% Minimal retained preamble for lane-x1 cards taken from
% sources/sigma-2023-005-on-asymptotically-locally-hyperbolic-metrics/pinned-original-source.tex
% (doi:10.3842/SIGMA.2023.005, SIGMA 19 (2023), article 005, CC BY 4.0).
%
% Retained verbatim from the source preamble: the theorem declarations
% under the per-section shared counter, the per-section equation numbering
% and the author operator/colour macros the retained spans use.
%
% Omitted and why: the \documentclass{sigma} lines (texlive carries no
% sigma.cls; per the section 8.8.6x precedent the assembler supplies the
% article baseline class) and the SIGMA front-matter commands, which belong
% to the journal layout and are not used by the retained spans.  The
% author's commented-out definition of \mycal is retained in active form:
% the publisher class that provided the macro is replaced, and the
% definition text is the author's own.
\usepackage{amsmath,amssymb,amsthm,mathrsfs,mathtools,graphicx}
\usepackage{xcolor}% the author's colour macros (red/blue/purple) are used in the retained spans
\usepackage{hyperref}% author bibliography entries use \href (8.18.2)

\numberwithin{equation}{section}

\newtheorem{Theorem}{Theorem}[section]
\newtheorem*{Theorem*}{Theorem}
\newtheorem{Corollary}[Theorem]{Corollary}
\newtheorem{Lemma}[Theorem]{Lemma}
\newtheorem{Proposition}[Theorem]{Proposition}
\newtheorem{Conjecture}[Theorem]{Conjecture}
 { \theoremstyle{definition}
\newtheorem{Definition}[Theorem]{Definition}
\newtheorem{Note}[Theorem]{Note}
\newtheorem{Example}[Theorem]{Example}
\newtheorem{Remark}[Theorem]{Remark} }

\newcommand{\blue}[1]{{\color{blue}#1}}
\newcommand{\hato}{\red{\omega}}
\newcommand{\hhi}{{\red{\hat h_{i}}}}
\newcommand{\hmv}{\red{\hat \omega}}
\newcommand{\horo}{\red{\mathfrak{h}}}
\newcommand{\hrho}{\hat\rho}
\newcommand{\hRi}{{\red{\hat R_{i}}}}
\newcommand{\mage}[1]{#1}%{\textcolor{magenta}}
\newcommand{\mcU}{{\mycal U}}
\newcommand{\mv}{\omega}%{\textcolor{magenta}{\omega}}
\newcommand{\mygb}{\red{b}}
\newcommand{\N}{\mathbb N}
\newcommand{\ominf}{\blue{\omega}}
\newcommand{\purple}[1]{{\color{purple}#1}}
\renewcommand{\qed}{}
\newcommand{\R}{\mathbb R}
\newcommand{\red}[1]{{\color{red}#1}}
\newcommand{\repsilon}{ {\varepsilon}}
\newcommand{\rhozero}{\red{\rho_0}}
\newcommand{\twodimN}{\purple{\mathbf{M}^2}}
\newcommand{\zMone}{\red{{\hat M}_{1,i}}}
\newcommand{\zMtwo}{\red{{\hat M}_{2,i}}}
\newcommand{\zomega}{\mathring{\omega}}%
\newcommand{\zui}{{\red{\hat u_{i}}}}
\newcommand{\mycal}{\cal}% the author's own (commented) definition, active here because the replaced publisher class provided the macro
